\section{Introduction}\label{sec:intro}

A systematic review begins with a search that returns thousands of records, and
every one of them has to be judged against pre-specified eligibility criteria
before a single study is read in full. This title and abstract screening stage
is consistently reported as the most time-consuming part of the review
process, and it is the stage where reviewer fatigue translates most directly
into inconsistent decisions \autocite{Tsafnat_2014}. The methodological
response has been to accept some degree of automation provided that the
automation is itself evaluated: the question is no longer whether machine
decisions may assist screening, but how their errors are measured and
reported.

Two families of tools have addressed the problem. Workflow systems such as
Rayyan, Covidence, DistillerSR and EPPI-Reviewer organise the screening work
but leave the judgement to the reviewer
\autocite{Ouzzani_2016,Marshall_2019}. Machine learning systems such as
ASReview, SWIFT-Review and the RCT classifiers derived from RobotSearch learn
from a small set of labelled records and then rank the remaining ones, so that
the reviewer reads the most promising records first
\autocite{van_de_Schoot_2021,Howard_2016}. Both families are useful, and both
have a structural limitation for the question this paper addresses: the
workflow systems produce no machine decision to evaluate, and the supervised
systems produce a ranking whose quality depends on the seed labels supplied by
the reviewer, which makes the result difficult to transfer across reviews.

Large language models change the shape of the problem. A general-purpose model
can be given the review protocol in the prompt and return a decision with a
short justification, without any task-specific training data. Early evidence
is encouraging but uneven. Khraisha and colleagues found that GPT-4 could
screen and extract data with performance that varied substantially by task and
by field \autocite{Khraisha_2024}; Gartlehner and colleagues reported a
similar pattern for data extraction \autocite{Gartlehner_2024}. At the same
time, Callaghan and colleagues argued that computer-assisted screening tools
are frequently released without the evaluation needed to justify their use in
evidence synthesis \autocite{Callaghan_2024}, and the journal has since
published explicit reporting expectations for manuscripts that evaluate
generative AI in systematic reviews and meta-analyses
\autocite{Farotimi_2025}.

That combination, promising capability plus a weak evidence base, defines the
gap this paper addresses. What a review team needs is not another model but an
instrument: a system in which the model can be replaced, in which every
automated decision is traceable to a record, a prompt and a model version, in
which cost is accounted per record rather than per project, and in which the
agreement between the model and a human reviewer can be measured on demand
rather than asserted in a paper. The absence of that instrument is, in our
reading, one reason why the literature on LLM screening is dominated by
one-off evaluations on curated datasets and thin on deployments inside real
reviews.

This paper describes \textsc{Sifter}, an open-source platform built to be that
instrument, and reports what a pilot deployment measured. The contribution is
threefold. First, we describe the design decisions that make LLM screening
auditable: a protocol object that is the single source for every prompt, a
three-way decision space in which \emph{doubt} is a first-class outcome
returned to the human reviewer, and a blind validation mode that hides the
model decision until the reviewer has recorded their own. Second, we report
the pilot evaluation: 14{,}219 records screened across four projects, the token
and monetary cost per record, and an agreement study on an educational corpus.
Third, we report the platform against the reporting expectations for
generative AI in evidence synthesis \autocite{Farotimi_2025} and state
explicitly which workload metrics, notably work saved over screening and
number needed to read, have not yet been measured.

The remainder of the paper is organised as follows.
Section~\ref{sec:related} situates the platform among existing review software
and screening automation. Section~\ref{sec:system} describes the architecture
and the screening pipeline. Section~\ref{sec:evaluation} reports the pilot
deployment. Section~\ref{sec:discussion} discusses what the numbers do and do
not establish, and Section~\ref{sec:conclusion} concludes.
